\chapter*{Cách học và kiến thức nền}
\addcontentsline{toc}{chapter}{Cách học và kiến thức nền}
\section*{Lộ trình học}
Hãy đọc phần nền này trước, rồi học Chương 1--3 theo thứ tự. Chương 4 dùng
tích vô hướng và trực chuẩn; Chương 5 dùng thêm bình phương tối thiểu.
Trong Chương 6, nhóm bài tính KDE chỉ cần định nghĩa, nhưng nhóm bandwidth
cần hiểu bias--variance. Với mỗi chứng minh, tự viết lại phép đổi biến và
kiểm tra điều kiện trước khi xem lời giải.

Mục tiêu của ``phi tham số'' là không cố định mật độ hoặc hàm hồi quy trong
một họ hữu hạn chiều như $\{N(\mu,\sigma^2)\}$. Ta vẫn có giả thiết:
mẫu độc lập, độ trơn, miền giá đỡ, moment và mật độ thiết kế. Càng ít
ràng buộc về hình dạng, càng cần dữ liệu để ước lượng cả một hàm.

\section*{Bảng ký hiệu}
\begin{longtable}{p{.23\textwidth}p{.70\textwidth}}
\toprule Ký hiệu & Ý nghĩa\\\midrule\endhead
$X_1,\ldots,X_n$ & Mẫu i.i.d. (độc lập, cùng phân phối), trừ khi ghi khác.\\
$F(x)$, $f(x)$ & CDF $\Pp(X\leq x)$ và mật độ, nếu phân phối tuyệt đối liên tục.\\
$\Fhat$, $\fhat$ & EDF và KDE; dấu mũ biểu thị đại lượng tính từ mẫu.\\
$h>0$, $D$ & Băng thông hoặc độ rộng ô; số chiều mô hình phép chiếu.\\
$p_j$, $\widehat p_j$ & Xác suất thật trong ô $B_j$; tần suất quan sát trong ô.\\
$K$, $\mu_j(K)$ & Kernel; moment $\int u^jK(u)\,du$.\\
$R(g)$ & $\int g(x)^2\,dx$; ở đây chữ $R$ trong $R(g)$ không phải ký hiệu rủi ro khác.\\
$\norm{g}_2^2$ & $\int g^2$ trên miền đang xét.\\
$\Bias$, $\Var$ & Chệch có dấu và phương sai lấy mẫu của ước lượng.\\
$\ISE$, $\MISE$ & Sai số bình phương tích phân của một mẫu; kỳ vọng của sai số đó.\\
$\AMSE$, $\AMISE$ & Xấp xỉ tiệm cận của MSE tại điểm và MISE toàn cục.\\
$\varphi$, $\Phi$ & Mật độ và CDF chuẩn tắc $N(0,1)$.\\
$r(x)$, $\sigma^2(x)$ & $\E(Y\mid X=x)$ và $\Var(Y\mid X=x)$.\\
\bottomrule
\end{longtable}

\section*{CDF, mật độ và xác suất}
Một CDF không giảm, liên tục phải và có giới hạn 0, 1 tại hai đầu trục số.
Nếu $F$ tuyệt đối liên tục thì $F(x)=\int_{-\infty}^x f(t)\,dt$,
$f\geq0$ và $\int f=1$. Khi đó
\[
\Pp(a<X\leq b)=F(b)-F(a)=\int_a^bf(t)\,dt.
\]
Mật độ là chiều cao, không phải xác suất tại điểm. Chẳng hạn mật độ đều
trên $[0,0.1]$ cao 10 nhưng diện tích là 1. Một phân phối rời rạc có CDF
nhưng không có mật độ theo độ đo Lebesgue; EDF luôn là phân phối rời rạc.

Với một hàm $a$ khả tích, $\int a\,dF$ là kỳ vọng $\E a(X)$.
Nếu có mật độ thì bằng $\int a(x)f(x)\,dx$; với EDF thì bằng
$n^{-1}\sum_i a(X_i)$. Ký hiệu $dF$ do đó dùng chung cho cả liên tục và rời rạc.

\section*{Biến chỉ báo và các quy tắc tính}
Đặt $I=\1_A$. Vì $I^2=I$,
\[
\E I=\Pp(A)=p,\qquad \Var I=\E I^2-(\E I)^2=p(1-p).
\]
Với hai sự kiện $A,B$,
$\Cov(\1_A,\1_B)=\Pp(A\cap B)-\Pp(A)\Pp(B)$.
Tuyến tính của kỳ vọng không đòi độc lập. Ngược lại,
\[
\Var\left(\sum_iZ_i\right)=\sum_i\Var Z_i+
2\sum_{i<j}\Cov(Z_i,Z_j).
\]
Chỉ khi các hiệp phương sai bằng 0 mới bỏ được tổng chéo.
Với i.i.d. có phương sai hữu hạn,
$\Var(n^{-1}\sum_iZ_i)=\Var(Z_1)/n$.

\section*{Chệch, phương sai và rủi ro}
Đặt $m=\E\widehat\theta$ và $\theta$ là giá trị thật.
\begin{align*}
\E(\widehat\theta-\theta)^2
&=\E[(\widehat\theta-m)+(m-\theta)]^2\\
&=\E(\widehat\theta-m)^2+(m-\theta)^2
 +2(m-\theta)\E(\widehat\theta-m)\\
&=\Var(\widehat\theta)+\Bias(\widehat\theta)^2.
\end{align*}
Hạng chéo bằng 0 vì $\E(\widehat\theta-m)=0$, không phải vì hai hạng
``độc lập''. Với hàm mật độ,
\[
\ISE=\int(\widehat f-f)^2,\qquad
\MISE=\E\ISE=\int\Bias(\widehat f(x))^2\,dx+
\int\Var(\widehat f(x))\,dx.
\]
Tonelli cho phép đổi kỳ vọng và tích phân của hàm không âm
$(\widehat f-f)^2$, kể cả khi tích phân vô hạn. Để tách hoặc đổi thứ tự
những tích phân có dấu, cần điều kiện khả tích tuyệt đối.
\intuition{Chệch đo vị trí trung bình của các ước lượng qua nhiều mẫu;
phương sai đo mức dao động quanh vị trí đó. Một đường ước lượng rất đẹp
trên một mẫu chưa chứng minh được rủi ro nhỏ.}

\section*{Hội tụ và ký hiệu tiệm cận}
\begin{itemize}
\item $Z_n\xrightarrow{a.s.}Z$: ngoài một sự kiện xác suất 0, dãy giá trị
  hội tụ. $Z_n\xrightarrow{P}Z$: với mọi $\epsilon>0$,
  $\Pp(|Z_n-Z|>\epsilon)\to0$.
\item $Z_n\Rightarrow Z$: hội tụ phân phối. CLT mô tả phân phối của
  sai số sau chuẩn hóa, không nói sai số chưa chuẩn hóa giữ nguyên độ lớn.
\item $a_n=O(b_n)$: $|a_n|\leq C|b_n|$ khi $n$ đủ lớn;
  $a_n=o(b_n)$: $a_n/b_n\to0$; $a_n\asymp b_n$: chặn trên và dưới
  bởi các hằng số dương.
\item $Z_n=O_P(b_n)$: $Z_n/b_n$ bị chặn theo xác suất;
  $Z_n=o_P(b_n)$: $Z_n/b_n\xrightarrow P0$.
\end{itemize}
LLN: nếu $Z_i$ i.i.d. và $\E|Z_1|<\infty$,
$\overline Z_n\xrightarrow{a.s.}\E Z_1$.
CLT: nếu thêm $0<\Var Z_1=\tau^2<\infty$,
\[
\sqrt n(\overline Z_n-\E Z_1)\Rightarrow N(0,\tau^2).
\]
Hội tụ điểm của $g_n(x)$ kiểm soát mỗi $x$ cố định; hội tụ đều kiểm soát
$\sup_x|g_n(x)-g(x)|$. Không tự đổi $x$ thành một điểm phụ thuộc mẫu
trong một kết quả chỉ có hội tụ điểm.

\section*{Taylor, delta method và phép đổi biến}
Nếu $g$ đủ trơn gần $x$,
\[
g(x+t)=g(x)+tg'(x)+\frac{t^2}2g''(x)+\cdots.
\]
Muốn tích phân phần dư trên toàn trục phải có chặn khả tích thích hợp.
Chỉ biết $g''$ liên tục cho phần dư $o(t^2)$, không tự cho $O(t^4)$.
Khi $u=(x-t)/h$, ta có $t=x-hu$, $dt=-h\,du$. Hai cận đảo thứ tự
nên dấu âm mất khi viết lại tích phân từ $-\infty$ đến $\infty$.

Delta method: nếu $\sqrt n(\widehat\theta-\theta)\Rightarrow N(0,\tau^2)$
và $g$ khả vi tại $\theta$, thì khi $g'(\theta)\ne0$,
\[
\sqrt n[g(\widehat\theta)-g(\theta)]
\Rightarrow N(0,[g'(\theta)]^2\tau^2).
\]
Nguồn gốc là khai triển $g(\widehat\theta)-g(\theta)
=g'(\theta)(\widehat\theta-\theta)+o_P(n^{-1/2})$.
Nếu đạo hàm bằng 0, phải xét bậc khai triển tiếp theo.

\section*{Tích vô hướng và trực chuẩn}
Trên $[0,1]$, đặt $\inner gq=\int_0^1gq$ và
$\norm g_2^2=\inner gg$. Nếu $\inner{\phi_j}{\phi_k}=\delta_{jk}$ thì
\[
\norm{\sum_{j=1}^Db_j\phi_j}_2^2
=\sum_{j,k}b_jb_k\inner{\phi_j}{\phi_k}=\sum_jb_j^2.
\]
Hạng chéo mất do tích phân các hàm trực giao, không do các hệ số ngẫu nhiên
độc lập. Đây là nền tảng của Chương 4.

\section*{Một phép tối ưu dùng xuyên suốt tài liệu}
Với $A,B,\beta>0$, xét $G(h)=Ah^{2\beta}+B/(nh)$.
\[
G'(h)=2\beta Ah^{2\beta-1}-\frac{B}{nh^2}=0
\quad\Longleftrightarrow\quad
h_*^{2\beta+1}=\frac{B}{2\beta An}.
\]
Vì $G(h)\to\infty$ khi $h\to0$ hoặc $h\to\infty$ và nghiệm dương duy
nhất, đây là cực tiểu toàn cục. Thế lại:
\[
h_*\asymp n^{-1/(2\beta+1)},\qquad
G(h_*)\asymp n^{-2\beta/(2\beta+1)}.
\]
Với $\beta=1$ nhận histogram, với $\beta=2$ nhận KDE bậc hai.
Đây là tối ưu một biểu thức xấp xỉ hoặc một chặn trên, không phải tự động
nghiệm tối ưu chính xác của rủi ro hữu hạn mẫu.
